\lesson{review}{让另一名智能体独立审查}
\diagram{p_review}
重要功能完成、合并多份改动或正式使用前，可安排独立审查。审查者直接读取要求与当前文件，自行核对结果；只看实现者的完成总结，容易沿用同一个误解。
\topic{将实现与审查分开}
\begin{enumerate}
\item \textbf{固定对象与材料。}记录版本、修改范围、需求和运行方法。优先用不继承实现对话的新会话或子智能体，只提供必要材料；复制原聊天历史的会话可能带入原有判断。
\item \textbf{按影响报告问题。}列出位置、触发条件、预期、实际和证据；优先处理正确性、安全与需求遗漏，风格偏好单列为可选。
\item \textbf{核对后修复，再复查。}实现者逐项验证意见，修复后交回当前版本。改动影响到的旧审查结论需重查；不为了回应每条建议而增加复杂度。
\end{enumerate}
\begin{prompt}{独立审查与复核}
材料：[原始要求、当前需求、版本、修改范围、运行方法]。若你参与实现，且工具支持，交给未参与且不继承实现对话的审查者；否则返回交接说明，可先自查但不称独立审查。材料不全先列缺口。\par
不改被审查实现。先读文件；需运行且会写入时，使用获准测试副本或独立工作目录及测试数据。禁止写入则列未运行项，不绕过限制；目录分开不代表外部数据已隔离。\par
核对要求、改动与证据，重点查错误结果、数据或权限问题、被删除或放宽的检查及旧功能受损；长期项目再查重复与复杂度。逐项给位置、触发条件、影响和证据，区分已验证问题与风险。\par
输出给实现者的清单；获准且有工具时交接，否则由用户转交。修复后核对当前版本再复查，汇总已解决、仍存在及未验证项；范围内未发现实质问题可说明，不为凑数修改。
\end{prompt}
\principle{独立视角也需要证据}{新审查者少受“自己刚做完”的判断影响，但仍可能犯相同错误。收到意见后，先核对触发步骤和证据，再交给实现者修复；修好后只复查受影响部分及项目要求的检查。多个 AI 同意不等于正确，拿不准的重要问题请有经验的人判断。}
